% capitulos/Capitulo2.tex — versión auditada e integrada
% Cambios principales: solución modal (Kundur §12.2), contexto WAMS (Sauer Cap. 10),
% Eckart–Young y Gavish–Donoho (Brunton & Kutz Cap. 1), VarPro (Brunton & Kutz §7.2, Askham–Kutz),
% corrección de notación (sigma_k para amortiguamiento), erratas y estructura.

\chapter{Marco Teórico}

\section*{Resumen del Capítulo}

Se estableció la cadena teórica completa que va de la ecuación de
oscilación al algoritmo de identificación: el modelo clásico multimáquina,
linealizado mediante los coeficientes de potencia sincronizante, produce un
sistema $\Delta\dot{\mathbf{x}} = \mathbf{A}\Delta\mathbf{x}$ cuya
respuesta libre es una suma de exponenciales amortiguadas
\eqref{eq:modal_response}; cualquier canal PMU hereda esa estructura
espectral \eqref{eq:measured_modal}, lo que justifica el modelo discreto
\eqref{eq:discrete} bajo la condición de Nyquist \eqref{eq:mapping}. El
algoritmo Multi-TK resuelve el problema de predicción lineal hacia atrás
multicanal \eqref{eq:multi_equation} mediante la pseudoinversa de
Moore--Penrose truncada \eqref{eq:pinv}, cuya optimalidad y efecto
regularizador quedan garantizados por el teorema de Eckart--Young--Mirsky
\eqref{eq:eckartyoung}; la separación geométrica de raíces en el plano $z$
distingue los modos físicos de los espurios, y los criterios de energía
\eqref{eq:energy}, umbral duro óptimo \eqref{eq:gavish_unknown} y energía
modal \eqref{eq:cum_mode_energy} gobiernan, respectivamente, el orden del
modelo y la jerarquización de modos dominantes. El refinamiento por
Proyección Variable \eqref{eq:varpro_projected}, inicializado con los polos
Multi-TK, suprime el sesgo de errores en variables de la etapa lineal y
optimiza directamente la métrica de reconstrucción; sus extensiones con
restricciones de caja y de consistencia de subespacio
\eqref{eq:P2}--\eqref{eq:P3} anclan el problema no convexo al subespacio de
señal identificado. Finalmente, las amplitudes complejas estimadas
\eqref{eq:amp_svd} se interpretan físicamente: la forma modal empírica es la
proyección del eigenvector derecho sobre los canales medidos
\eqref{eq:shape_link}, el índice $\hat{p}_{ik}$ cuantifica observabilidad
relativa --- distinto, por identificabilidad, del factor de participación
de estados \eqref{eq:participation_state} --- y los cosenos directores
\eqref{eq:direction_cosine} junto con el agrupamiento espectral por
eigengap \eqref{eq:eigengap} particionan el sistema en grupos coherentes de
generadores sin parámetros fijados a priori. Este marco, validado contra la
literatura de identificación basada en sincrofasores
\cite{jiang2021,li2018cwt,li2020era}, constituye el fundamento del
algoritmo Multi-TK + OptVarPro desarrollado en los capítulos siguientes.


% ---------------- INTRODUCCION DEL CAPITULO ----------------
\section*{Introducción del Capítulo}

Este capítulo construye, de manera deductiva, el puente entre la física de
las oscilaciones electromecánicas en sistemas eléctricos de potencia y los
algoritmos de identificación modal basados en mediciones sincrofasoriales
que sustentan esta tesis. La exposición sigue una cadena de implicaciones:
se parte de las ecuaciones de oscilación del generador síncrono bajo el
modelo clásico \cite{kundur1994}, se linealiza alrededor del punto de
operación para obtener el modelo de estados cuyo espectro define los modos
electromecánicos, y se demuestra --- mediante la transformación modal y la
respuesta libre del sistema \cite{kundur1994} --- que toda señal registrada
por una unidad de medición fasorial en régimen de \emph{ring-down} obedece
exactamente un modelo de suma de exponenciales complejas amortiguadas. Esta
derivación convierte el modelo de señal de los métodos de identificación en
una consecuencia de la física del sistema, y no en un postulado, y delimita
con precisión su dominio de validez: la banda electromecánica
($0.1$--$3$~Hz) observada con tasas de reporte PMU de $30$~Hz o superiores
\cite{sauer2017}.

Sobre esa base, el capítulo desarrolla la maquinaria matemática de
estimación: la formulación discreta del modelo exponencial y su mapeo
biunívoco al plano $z$; la predicción lineal hacia atrás multicanal del
algoritmo Tufts--Kumaresan, cuyo mecanismo de separación señal--ruido en el
plano complejo se analiza en detalle \cite{tufts1982}; el truncamiento de
la descomposición en valores singulares como filtro óptimo en el sentido de
Eckart--Young--Mirsky, junto con los criterios de selección de orden por
energía acumulada y por umbral duro óptimo \cite{brunton2022,gavish2014};
la estimación de amplitudes complejas mediante sistemas de Vandermonde
numéricamente estables; y el refinamiento por mínimos cuadrados no lineales
separables mediante Proyección Variable, que mitiga el sesgo estructural de
los estimadores lineales \cite{askham2018}. La parte final traduce los
parámetros estimados en \emph{patrones dinámicos} interpretables para la
operación del sistema: formas modales y su vínculo formal con los
eigenvectores derechos, índices de participación basados en mediciones y su
distinción rigurosa frente a los factores de participación del análisis de
pequeña señal, y la identificación de grupos coherentes de generadores
mediante cosenos directores y agrupamiento espectral con selección
automática del número de grupos \cite{jiang2021,li2018cwt,li2020era}.

El capítulo incluye, además, el diseño de los estudios experimentales que
validan las decisiones metodológicas: la selección del orden del predictor
$L$ mediante diagramas de estabilización, la del orden del modelo $M$
contrastando el criterio de energía con el umbral de Gavish--Donoho, y las
formulaciones de optimización con restricciones de subespacio que acoplan
la etapa lineal y la no lineal del método propuesto.



\section{Dinámica de Oscilaciones Electromecánicas}
% ===========================================================

\subsection{Modelado del Generador Síncrono}
\subsubsection{Ecuaciones de Swing del Generador Síncrono}

La dinámica de cada generador $i$ en un sistema de $n_g$ máquinas ($i=1,\ldots,n_g$) se describe mediante las \textbf{ecuaciones de swing}:

\begin{align}
    \dot{\delta}_i  &= \omega_i - \omega_s \label{eq:swing1} \\[4pt]
    \dot{\omega}_i  &= \frac{\omega_s}{2H_i}
                       \Bigl(P_{m,i} - P_{e,i}(\boldsymbol{\delta}) - D_i(\omega_i - \omega_s)\Bigr)
                       \label{eq:swing2}
\end{align}

Las ecuaciones \eqref{eq:swing1}--\eqref{eq:swing2} corresponden al
\textbf{modelo clásico} del generador síncrono \cite{kundur1994}: la máquina
se representa como una fuente de voltaje interna de magnitud constante
$E'_i$ detrás de la reactancia transitoria $x'_{d,i}$, con potencia mecánica
$P_{m,i}$ constante en la ventana de análisis y con todos los efectos de
amortiguamiento (devanados amortiguadores, cargas sensibles a la frecuencia, etc.)
agregados en el coeficiente $D_i$. Este nivel de modelado es suficiente para
caracterizar la dinámica electromecánica de interés, pues los modos de
oscilación de nuestro interes  entre rotores quedan determinados por el intercambio de energía
cinética y potencia sincronizante entre máquinas \cite{kundur1994,sauer2017}.

% -----------------------------------------------------------

Parámetros (todas las potencias en por unidad sobre una base común
$S_{\text{base}}$):

\begin{itemize}
    \item $\delta_i$: ángulo del rotor respecto a la referencia síncrona (rad);
    \item $\omega_i$: velocidad angular del rotor (rad/s);
    \item $H_i$: constante de inercia, $H_i = \tfrac{1}{2}J_i\omega_s^2/S_{\text{base}}$ (s);
    \item $D_i$: coeficiente de amortiguamiento (p.u.\ de par por
          p.u.\ de desviación de velocidad, adimensional);
    \item $P_{m,i}$: potencia mecánica de entrada (p.u.);
    \item $P_{e,i}(\boldsymbol{\delta})$: potencia eléctrica entregada ---
          función no lineal de los ángulos de rotor
          $\boldsymbol{\delta} = [\delta_1, \ldots, \delta_{n_g}]^\top$ a
          través de la red a nodos de generación (p.u.);
    \item $\omega_s = 2\pi f_s$: velocidad síncrona de referencia (rad/s).
\end{itemize}

% -----------------------------------------------------------

\subsubsection{Linealización alrededor del Punto de Operación}

Para analizar la estabilidad de pequeña señal, se linealiza el sistema alrededor del punto de equilibrio
$(\boldsymbol{\delta}^0, \boldsymbol{\omega}^0)$; en régimen permanente se cumple ($\dot{\delta}_i = 0$, $\dot{\omega}_i = 0$):
\begin{equation}
    P_{m,i} = P_{e,i}(\boldsymbol{\delta}^0), \qquad \omega_i^0 = \omega_s
\end{equation}

Se definen las variables de estado de perturbación (desviaciones):
\begin{equation}
    \Delta\boldsymbol{\delta} = \boldsymbol{\delta} - \boldsymbol{\delta}^0, \qquad
    \Delta\boldsymbol{\omega} = \boldsymbol{\omega} - \boldsymbol{\omega}^0
\end{equation}

% -----------------------------------------------------------

Linealizando la primera ecuación de oscilación:
\begin{equation}
    \Delta\dot{\boldsymbol{\delta}} = \Delta\boldsymbol{\omega}
\end{equation}

Expandiendo la serie de Taylor de primer orden de $P_{e,i}(\boldsymbol{\delta})$ alrededor de $\boldsymbol{\delta}^0$ y usando la condición de equilibrio se linealiza la segunda ecuación de oscilación:
\begin{equation}
    \Delta\dot{\omega}_i = \frac{\omega_s}{2H_i}\Bigl(-\sum_{j=1}^{n_g} \frac{\partial P_{e,i}}{\partial \delta_j}\Big|_0 \Delta\delta_j - D_i \Delta\omega_i\Bigr)
\end{equation}

Las derivadas parciales evaluadas en el punto de operación definen los
\textbf{coeficientes de potencia sincronizante} \cite{kundur1994}:
\begin{equation}
    K_{ij} \;=\; \left.\frac{\partial P_{e,i}}{\partial \delta_j}\right|_{\boldsymbol{\delta}^0},
    \qquad i,j = 1,\ldots,n_g,
    \label{eq:Kij}
\end{equation}
que representan el acoplamiento entre las
máquinas $i$ y $j$ a través de la red y representa el cambio en la potencia eléctrica para un cambio en el punto de equilibr. Agrupando las desviaciones en el
vector de estados $\Delta\mathbf{x} =
[\Delta\boldsymbol{\delta}^\top \;\; \Delta\boldsymbol{\omega}^\top]^\top
\in \mathbb{R}^{2n_g}$, el sistema linealizado adquiere la forma canónica
de espacio de estados en lazo abierto (movimiento libre):
\begin{equation}
    \Delta\dot{\mathbf{x}} \;=\; \mathbf{A}\,\Delta\mathbf{x},
    \qquad
    \mathbf{A} \;=\;
    \begin{bmatrix}
        \mathbf{0} & \mathbf{I}_{n_g} \\[4pt]
        -\,\omega_s\,\mathbf{M}^{-1}\mathbf{K} & -\,\omega_s\,\mathbf{M}^{-1}\mathbf{D}
    \end{bmatrix}
    \in \mathbb{R}^{2n_g \times 2n_g},
    \label{eq:statespace}
\end{equation}
donde $\mathbf{M} = \operatorname{diag}(2H_1,\ldots,2H_{n_g})$,
$\mathbf{D} = \operatorname{diag}(D_1,\ldots,D_{n_g})$ y
$\mathbf{K} = [K_{ij}] \in \mathbb{R}^{n_g\times n_g}$ es la matriz de
coeficientes de potencia sincronizante. La estabilidad de pequeña señal del
punto de operación queda gobernada por los eigenvalores de $\mathbf{A}$,
raíces de la ecuación característica
$\det(s\mathbf{I}-\mathbf{A}) = 0$ \cite{kundur1994}.

% -----------------------------------------------------------

\subsubsection{Solución Modal: del Espacio de Estados a la Suma de
Exponenciales Amortiguadas}

Sea $\{\lambda_k\}_{k=1}^{2n_g}$ el espectro de $\mathbf{A}$, con
eigenvectores derechos $\boldsymbol{\phi}_k$ e izquierdos
$\boldsymbol{\psi}_k$ que satisfacen
\begin{equation}
    \mathbf{A}\boldsymbol{\phi}_k = \lambda_k \boldsymbol{\phi}_k,
    \qquad
    \boldsymbol{\psi}_k^\top \mathbf{A} = \lambda_k \boldsymbol{\psi}_k^\top,
    \qquad
    \boldsymbol{\psi}_j^\top\boldsymbol{\phi}_k = \delta_{jk}.
    \label{eq:eigen}
\end{equation}
La transformación $\Delta\mathbf{x} = \boldsymbol{\Phi}\mathbf{z}$,
con $\boldsymbol{\Phi}=[\boldsymbol{\phi}_1 \cdots \boldsymbol{\phi}_{2n_g}]$,
desacopla \eqref{eq:statespace} en $2n_g$ ecuaciones escalares
$\dot{z}_k = \lambda_k z_k$, de modo que la respuesta libre ante una
perturbación inicial $\Delta\mathbf{x}(0)$ es \cite{kundur1994}
\begin{equation}
    \Delta\mathbf{x}(t)
    \;=\;
    \sum_{k=1}^{2n_g} \boldsymbol{\phi}_k\,
    \underbrace{\bigl(\boldsymbol{\psi}_k^\top \Delta\mathbf{x}(0)\bigr)}_{c_k}\,
    e^{\lambda_k t}.
    \label{eq:modal_response}
\end{equation}

Dado que $\mathbf{A}$ es real, los eigenvalores oscilatorios aparecen en
pares complejos conjugados $\lambda_k = \sigma_k \pm j\omega_k$; cada par
constituye un \textbf{modo electromecánico de oscilación} con frecuencia
$f_k = \omega_k/2\pi$ y razón de amortiguamiento
\begin{equation}
    \zeta_k \;=\; \frac{-\sigma_k}{\sqrt{\sigma_k^2 + \omega_k^2}}.
    \label{eq:damping_ratio}
\end{equation}
Cualquier señal medida $x_i(t)$ que sea combinación lineal de los estados
--- es el caso de desviaciones de velocidad de rotor, frecuencia de bus,
potencia activa o ángulos de tensión registrados por un PMU, expresables
como $x_i = \mathbf{c}_i^\top\Delta\mathbf{x}$ --- tiene la misma
estructura espectral:
\begin{equation}
    x_i(t) \;=\; \sum_{k=1}^{2n_g}
    \underbrace{\bigl(\mathbf{c}_i^\top\boldsymbol{\phi}_k\bigr)\,c_k}_{a_{i,k}}\,
    e^{\lambda_k t},
    \label{eq:measured_modal}
\end{equation}
lo que justifica, desde la física del sistema, el modelo de suma de
exponenciales amortiguadas que adoptan los algoritmos de identificación
modal de las secciones siguientes. La amplitud compleja $a_{i,k}$ acopla la
\textit{observabilidad} del modo $k$ en el canal $i$
(a través de $\mathbf{c}_i^\top\boldsymbol{\phi}_k$, es decir, de la forma
modal) con la \textit{excitación} del modo por la perturbación
(a través de $c_k$) \cite{kundur1994}.

% -----------------------------------------------------------

\subsubsection{Modos Electromecánicos y su Observación mediante
Sincrofasores (WAMS)}

Los modos oscilatorios de \eqref{eq:modal_response} se clasifican según su
extensión geográfica y banda de frecuencia \cite{kundur1994,sauer2017}:
\textbf{modos locales} (una planta contra el resto del sistema,
aproximadamente $1$--$3$~Hz) y \textbf{modos interárea} (grupos coherentes
de generadores oscilando entre sí, por debajo de $1$~Hz, típicamente
$0.25$--$1$~Hz). Los modos interárea, de bajo amortiguamiento
($\zeta_k < 3\%$) y difícil diagnóstico, limitan las transferencias de
potencia y constituyen el objeto principal de la identificación modal
\cite{sauer2017}.

Las unidades de medición fasorial (PMU) de un sistema de monitoreo de área
amplia (WAMS) entregan fasores de tensión y corriente, frecuencia y
potencias sincronizados vía GPS, con tasas de reporte de $30$ a $120$
muestras por segundo \cite{sauer2017}. Puesto que la banda electromecánica
completa se encuentra por debajo de $3$~Hz, una tasa de reporte
$F_s \ge 30$~Hz satisface holgadamente el criterio de Nyquist
($F_s > 2 f_k$) y dota a las señales sincrofasoriales de la resolución
temporal necesaria para estimar tanto $\omega_k$ como $\sigma_k$
\cite{sauer2017}. En contraste, los sistemas SCADA convencionales, con
periodos de muestreo del orden de segundos, no resuelven esta dinámica.

Tras una gran perturbación, la respuesta inmediata del sistema es
fuertemente no lineal; sin embargo, una vez que las trayectorias se
aproximan al nuevo punto de equilibrio, las oscilaciones decrecientes
registradas por los PMU --- el régimen de \textbf{ring-down} --- se modelan
con precisión como la respuesta libre de un sistema lineal no forzado
\cite{sauer2017}. Es en este régimen donde el modelo
\eqref{eq:measured_modal} es válido y donde operan los algoritmos de
predicción lineal desarrollados a continuación.

\section{Formulación Discreta y Predicción Lineal}
% ===========================================================

\subsection{Del Tiempo Continuo al Dominio Muestreado}

Muestreando cada señal $x_i(t)$ de \eqref{eq:measured_modal} con periodo
$\Delta t = 1/F_s$, definimos $x_i[n] = x_i(n\Delta t)$ para
$n = 0,1,\dots,N-1$. El modelo continuo de suma de exponenciales
amortiguadas se transforma exactamente en
\begin{equation}
    x_i[n] = \sum_{k=1}^{M} a_{i,k}\, z_k^{\,n}, \qquad
    z_k = e^{\lambda_k \Delta t} = e^{(\sigma_k + j\omega_k)\Delta t},
    \label{eq:discrete}
\end{equation}
donde $M \le 2n_g$ es el número de modos efectivamente observables en la
ventana de datos. Puesto que las señales medidas son reales, los polos
$z_k$ y las amplitudes $a_{i,k}$ aparecen en pares complejos conjugados,
por lo que $M$ es par para modos puramente oscilatorios. Los polos en
tiempo discreto contienen toda la información modal:

\begin{align}
    z_k &= e^{(\sigma_k + j\omega_k)\Delta t}, \\[4pt]
    \lambda_k = \sigma_k + j\omega_k &= \frac{\ln z_k}{\Delta t},
    \qquad |\omega_k|\,\Delta t < \pi .
    \label{eq:mapping}
\end{align}

La condición $|\omega_k|\Delta t < \pi$ es el criterio de Nyquist que
garantiza la unicidad de la rama principal del logaritmo complejo en
\eqref{eq:mapping}; las tasas de reporte PMU ($F_s \ge 30$~Hz) la
satisfacen con amplio margen para la banda electromecánica
($f_k \le 3$~Hz) \cite{sauer2017}. A partir de las raíces $z_k$ del
polinomio de predicción se recuperan el amortiguamiento
$\sigma_k = \operatorname{Re}(\ln z_k)/\Delta t$, la frecuencia
$\omega_k = \operatorname{Im}(\ln z_k)/\Delta t$ y la razón de
amortiguamiento $\zeta_k$ de \eqref{eq:damping_ratio}.

\section{Algoritmo Tufts-Kumaresan Multivariable (Multi-TK)}
% ===========================================================

\subsection{Modelo de Señal y Predicción Lineal hacia Atrás}

Para $P$ señales medidas $\{x_i[n]\}_{i=1}^P$, cada una de $N$ muestras ($n=0,\dots,N-1$), el modelo de señal con ruido de medición $w_i[n]$ es
\[
x_i[n] = \sum_{k=1}^{M} a_{i,k}\, z_k^n + w_i[n], \qquad z_k = e^{(\sigma_k + j\omega_k)\Delta t}.
\]

Se elige un orden de predictor $L$ tal que $M \le L \le N-M$. Para cada canal $i$ se plantea el problema de \textbf{predicción lineal hacia atrás}:
\begin{equation}
\mathbf{A}_i\,\mathbf{b} = -\mathbf{h}_i,
\label{eq:backward_pred}
\end{equation}
donde $\mathbf{b}=[b(1),b(2),\dots,b(L)]^\top$ es el vector de coeficientes del polinomio característico.

La elección de la predicción \textit{hacia atrás} sobre datos conjugados
no es accidental: constituye el mecanismo de discriminación señal--ruido
del método de Tufts--Kumaresan \cite{tufts1982}. Al elegir un orden
sobredimensionado $L > M$, el polinomio característico
\begin{equation}
    B(z) \;=\; 1 + \sum_{\ell=1}^{L} b(\ell)\, z^{-\ell}
    \label{eq:charpoly}
\end{equation}
posee $L$ raíces, de las cuales sólo $M$ corresponden a modos físicos. En
la formulación hacia atrás, las raíces de señal se ubican en los
recíprocos conjugados de los polos, $1/z_k^*$: para modos amortiguados
($|z_k|<1$) éstas quedan \textbf{fuera} del círculo unitario, mientras que
las $L-M$ raíces extrañas inducidas por el ruido tienden a permanecer
\textbf{dentro} de él \cite{tufts1982}. Esta separación geométrica en el
plano $z$ permite identificar los modos físicos sin ambigüedad, propiedad
de la que carece la predicción hacia adelante del método de Prony clásico
\cite{sauer2017,hauer1990}, cuya estructura Hankel es por lo demás análoga
a \eqref{eq:Ai}.

La matriz de Hankel conjugada para la $i$-ésima señal (el asterisco denota
conjugación compleja) se construye como:
\begin{equation}
\mathbf{A}_i =
\begin{bmatrix}
x_i^*(1)   & x_i^*(2)   & \cdots & x_i^*(L) \\
x_i^*(2)   & x_i^*(3)   & \cdots & x_i^*(L+1) \\
\vdots     & \vdots     & \ddots & \vdots   \\
x_i^*(N-L) & x_i^*(N-L+1) & \cdots & x_i^*(N-1)
\end{bmatrix}
\in \mathbb{C}^{(N-L)\times L},
\label{eq:Ai}
\end{equation}
y el vector objetivo asociado es
\begin{equation}
\mathbf{h}_i =
\begin{bmatrix}
x_i^*(0) \\ x_i^*(1) \\ \vdots \\ x_i^*(N-L-1)
\end{bmatrix}
\in \mathbb{C}^{N-L}.
\label{eq:hi}
\end{equation}

\subsection{Sistema Global Multicanal}

Se apilan verticalmente las matrices y los vectores de todos los canales:
\begin{equation}
\mathbf{A}_{\text{multi}} =
\begin{bmatrix} \mathbf{A}_1 \\ \mathbf{A}_2 \\ \vdots \\ \mathbf{A}_P \end{bmatrix}
\in \mathbb{C}^{P(N-L)\times L},
\qquad
\mathbf{h}_{\text{multi}} =
\begin{bmatrix} \mathbf{h}_1 \\ \mathbf{h}_2 \\ \vdots \\ \mathbf{h}_P \end{bmatrix}
\in \mathbb{C}^{P(N-L)}.
\label{eq:multi_system}
\end{equation}

El sistema completo de predicción lineal es:
\begin{equation}
\boxed{\mathbf{A}_{\text{multi}}\,\mathbf{b} = -\mathbf{h}_{\text{multi}}}.
\label{eq:multi_equation}
\end{equation}

\subsection{Descomposición SVD y Selección del Orden del Modelo}

Se aplica la \textbf{descomposición en valores singulares} a $\mathbf{A}_{\text{multi}}$:
\begin{equation}
\mathbf{A}_{\text{multi}} = \mathbf{U}\,\boldsymbol{\Sigma}\,\mathbf{V}^H,
\label{eq:svd}
\end{equation}
con $\boldsymbol{\Sigma} = \operatorname{diag}(\sigma_1,\sigma_2,\dots,\sigma_r)$, $\sigma_1 \ge \sigma_2 \ge \dots \ge \sigma_r > 0$
(en esta sección, $\sigma_k$ denota los valores singulares de
$\mathbf{A}_{\text{multi}}$, no el amortiguamiento modal).

La optimalidad del truncamiento de la SVD está garantizada por el teorema
de Eckart--Young--Mirsky \cite{brunton2022}: la mejor aproximación de
rango $M$ a una matriz, en el sentido de mínimos cuadrados, es
precisamente su SVD truncada.

Sea $\mathbf{A}_{\mathrm{multi}} = \mathbf{U}\boldsymbol{\Sigma}\mathbf{V}^H$.
Entonces
\begin{equation}
    \mathbf{U}_s\boldsymbol{\Sigma}_s\mathbf{V}_s^H
    \;=\;
    \operatorname*{arg\,min}_{\operatorname{rank}(\mathbf{B}) \,\le\, M}
    \bigl\lVert \mathbf{A}_{\mathrm{multi}} - \mathbf{B} \bigr\rVert_F,
    \label{eq:eckartyoung}
\end{equation}
y el error de aproximación queda exactamente cuantificado por la energía
de los valores singulares descartados:
\begin{equation}
    \bigl\lVert \mathbf{A}_{\mathrm{multi}}
    - \mathbf{U}_s\boldsymbol{\Sigma}_s\mathbf{V}_s^H \bigr\rVert_F^2
    \;=\; \sum_{k=M+1}^{r} \sigma_k^2 .
    \label{eq:trunc_error}
\end{equation}

Bajo el modelo aditivo
$\mathbf{A}_{\mathrm{multi}} = \mathbf{A}_{\mathrm{señal}} +
\mathbf{A}_{\mathrm{ruido}}$, donde
$\operatorname{rank}(\mathbf{A}_{\mathrm{señal}}) = M$ por construcción del
modelo \eqref{eq:discrete}, los primeros $M$ valores singulares capturan el
subespacio de señal y los restantes la contribución del ruido de medición;
el truncamiento de rango $M$ actúa entonces como un filtro óptimo en el
sentido de \eqref{eq:eckartyoung} \cite{brunton2022,tufts1982}.

El número de modos $M$ se selecciona mediante el criterio de \textbf{energía acumulada}:
\begin{equation}
E_{s,M} = \frac{\displaystyle\sum_{k=1}^{M} \sigma_k^2}{\displaystyle\sum_{k=1}^{r} \sigma_k^2} \ge E_0,
\label{eq:energy}
\end{equation}
donde típicamente $E_0 = 97\%$ (o $99.9\%$ según el nivel de ruido).

El umbral de energía $E_0$ es una regla práctica cuyo valor óptimo depende
del nivel de ruido. Una alternativa con garantía estadística es el
\textbf{umbral duro óptimo} de Gavish--Donoho \cite{gavish2014,brunton2022},
derivado bajo el modelo
$\mathbf{A}_{\mathrm{multi}} = \mathbf{A}_{\mathrm{señal}} +
\gamma\,\mathbf{A}_{\mathrm{ruido}}$, con entradas de ruido i.i.d.\
gaussianas de media cero y varianza unitaria. Para una matriz rectangular
de dimensiones $n\times m$ ($\beta = m/n \le 1$) y nivel de ruido $\gamma$
conocido, se retienen únicamente los valores singulares que superan
\begin{equation}
    \tau \;=\; \lambda(\beta)\,\sqrt{n}\,\gamma,
    \qquad
    \lambda(\beta) = \left( 2(\beta+1) +
    \frac{8\beta}{(\beta+1) + \sqrt{\beta^2 + 14\beta + 1}} \right)^{1/2},
    \label{eq:gavish_known}
\end{equation}
expresión que se reduce a $\tau = (4/\sqrt{3})\sqrt{n}\,\gamma$ en el caso
cuadrado. Cuando $\gamma$ es desconocido --- la situación habitual con
mediciones sincrofasoriales --- el umbral se estima a partir de la mediana
de los valores singulares, $\sigma_{\mathrm{med}}$:
\begin{equation}
    \tau \;=\; \omega(\beta)\,\sigma_{\mathrm{med}},
    \label{eq:gavish_unknown}
\end{equation}
donde $\omega(\beta)$ se calcula numéricamente \cite{gavish2014}. En este
trabajo, el criterio de energía \eqref{eq:energy} se contrasta con
\eqref{eq:gavish_unknown} como verificación cruzada del orden $M$.

\subsection{Truncamiento de la SVD y Solución del Sistema}

Se retiene el subespacio de señal de rango $M$:
\begin{equation}
\mathbf{A}_{\text{multi}} \approx \mathbf{U}_s \boldsymbol{\Sigma}_s \mathbf{V}_s^H,
\quad
\mathbf{U}_s \in \mathbb{C}^{P(N-L)\times M},\;
\boldsymbol{\Sigma}_s \in \mathbb{R}^{M\times M},\;
\mathbf{V}_s \in \mathbb{C}^{L\times M}.
\label{eq:truncsvd}
\end{equation}

La solución regularizada del sistema es:
\begin{equation}
\boxed{\mathbf{b} = -\mathbf{V}_s\,\boldsymbol{\Sigma}_s^{-1}\,\mathbf{U}_s^H\,\mathbf{h}_{\text{multi}}}.
\label{eq:solution_final}
\end{equation}

La ecuación \eqref{eq:solution_final} es la solución de mínimos cuadrados
mediante la \textbf{pseudoinversa de Moore--Penrose truncada}
\cite{brunton2022}:
\begin{equation}
    \mathbf{b}
    \;=\;
    -\,\mathbf{A}_{\mathrm{multi}}^{+(M)}\,\mathbf{h}_{\mathrm{multi}},
    \qquad
    \mathbf{A}_{\mathrm{multi}}^{+(M)}
    \;=\;
    \mathbf{V}_s\,\boldsymbol{\Sigma}_s^{-1}\,\mathbf{U}_s^H .
    \label{eq:pinv}
\end{equation}
Esta solución posee una doble optimalidad \cite{brunton2022}: minimiza el
residuo $\lVert \mathbf{A}_{\mathrm{multi}}\mathbf{b} +
\mathbf{h}_{\mathrm{multi}} \rVert_2$ restringido al subespacio de señal de
rango $M$ y, entre todas las soluciones que lo logran, es la de mínima
norma $\lVert\mathbf{b}\rVert_2$. El truncamiento cumple además una función
de regularización numérica: en la pseudoinversa completa, los valores
singulares pequeños asociados al ruido aparecen invertidos
($\sigma_k^{-1}$ grande) y amplifican catastróficamente la componente de
ruido de $\mathbf{h}_{\mathrm{multi}}$; al retener sólo
$\boldsymbol{\Sigma}_s$, dicha amplificación se elimina y la varianza del
estimador de $\mathbf{b}$ se reduce drásticamente
\cite{tufts1982,brunton2022}. Para la longitud del predictor se recomienda
$L \approx \tfrac{3N}{4}$, valor que maximiza el desempeño del estimador en
presencia de ruido \cite{tufts1982}.

\section{Patrones Dinámicos: Formas Modales y Factores de Participación}
% ===========================================================

\subsection{Estimación de Amplitudes Complejas}

Conocidos los polos $\{z_k\}_{k=1}^{M}$, para cada señal $i$ se resuelve:

\begin{equation}
\underbrace{
\begin{bmatrix}
    z_1^0 & z_2^0 & \cdots & z_M^0 \\
    z_1^1 & z_2^1 & \cdots & z_M^1 \\
    \vdots &      & \ddots & \vdots \\
    z_1^{N-1} & z_2^{N-1} & \cdots & z_M^{N-1}
\end{bmatrix}
}_{\boldsymbol{\Psi}}
\begin{bmatrix} a_{i,1} \\ a_{i,2} \\ \vdots \\ a_{i,M} \end{bmatrix}
= \mathbf{x}_i
\end{equation}

Mediante mínimos cuadrados:
\begin{equation}
\mathbf{a}_i = \boldsymbol{\Psi}^+\,\mathbf{x}_i
= \bigl(\boldsymbol{\Psi}^H\boldsymbol{\Psi}\bigr)^{-1}\boldsymbol{\Psi}^H\mathbf{x}_i
\end{equation}

La matriz $\boldsymbol{\Psi}\in\mathbb{C}^{N\times M}$ es una matriz de
\textbf{Vandermonde} construida sobre los polos estimados. En la práctica,
la solución vía ecuaciones normales
$(\boldsymbol{\Psi}^H\boldsymbol{\Psi})^{-1}\boldsymbol{\Psi}^H$ debe
evitarse: cuando el sistema presenta modos de frecuencias cercanas ---
situación típica de los modos interárea --- los nodos $z_k$ se aglomeran y
$\boldsymbol{\Psi}$ se vuelve mal condicionada, y el producto
$\boldsymbol{\Psi}^H\boldsymbol{\Psi}$ eleva al cuadrado su número de
condición, $\kappa(\boldsymbol{\Psi}^H\boldsymbol{\Psi}) =
\kappa(\boldsymbol{\Psi})^2$. La alternativa numéricamente estable es la
pseudoinversa calculada mediante la SVD de $\boldsymbol{\Psi}$
\cite{brunton2022}:
\begin{equation}
    \mathbf{a}_i \;=\;
    \mathbf{V}_{\Psi}\,\boldsymbol{\Sigma}_{\Psi}^{-1}\,
    \mathbf{U}_{\Psi}^{H}\,\mathbf{x}_i,
    \qquad
    \boldsymbol{\Psi} =
    \mathbf{U}_{\Psi}\boldsymbol{\Sigma}_{\Psi}\mathbf{V}_{\Psi}^{H}.
    \label{eq:amp_svd}
\end{equation}

\subsection{Refinamiento por Mínimos Cuadrados No Lineales: Proyección
Variable}

Los estimadores de la familia de predicción lineal comparten una limitación
estructural: el ruido de medición contamina simultáneamente la matriz de
regresores $\mathbf{A}_{\mathrm{multi}}$ y el vector objetivo
$\mathbf{h}_{\mathrm{multi}}$ --- un problema de \textit{errores en
variables} --- por lo que la solución de mínimos cuadrados lineal produce
estimaciones de los polos con \textbf{sesgo} que el truncamiento de la SVD
reduce pero no elimina \cite{brunton2022}. La formulación que ataca el
sesgo de raíz consiste en resolver directamente el problema de ajuste
exponencial como un problema de mínimos cuadrados \textbf{no lineal} en los
parámetros modales \cite{askham2018,brunton2022}. Organizando las $P$
señales en la matriz de datos
$\mathbf{X} = [\mathbf{x}_1\;\cdots\;\mathbf{x}_P] \in \mathbb{C}^{N\times P}$
y las amplitudes en $\mathbf{a} \in \mathbb{C}^{M\times P}$, el problema es
\begin{equation}
    \min_{\boldsymbol{\lambda},\,\mathbf{a}}\;
    \bigl\lVert \mathbf{X} -
    \boldsymbol{\Psi}(\boldsymbol{\lambda})\,\mathbf{a} \bigr\rVert_F^2,
    \qquad
    [\boldsymbol{\Psi}(\boldsymbol{\lambda})]_{n,k} = e^{\lambda_k n \Delta t},
    \label{eq:varpro_problem}
\end{equation}

donde $\boldsymbol{\lambda} = [\lambda_1,\ldots,\lambda_M]^\top$. El
problema \eqref{eq:varpro_problem} es \textbf{separable}: para
$\boldsymbol{\lambda}$ fijo, las amplitudes óptimas tienen solución cerrada
$\hat{\mathbf{a}}(\boldsymbol{\lambda}) =
\boldsymbol{\Psi}(\boldsymbol{\lambda})^{+}\mathbf{X}$. El método de
\textbf{Proyección Variable} (VarPro) explota esta estructura eliminando
las variables lineales y optimizando únicamente sobre las no lineales
\cite{askham2018}:
\begin{equation}
    \min_{\boldsymbol{\lambda}}\;
    \Bigl\lVert
    \bigl(\mathbf{I} - \boldsymbol{\Psi}(\boldsymbol{\lambda})
    \boldsymbol{\Psi}(\boldsymbol{\lambda})^{+}\bigr)\,\mathbf{X}
    \Bigr\rVert_F^2,
    \label{eq:varpro_projected}
\end{equation}
es decir, se minimiza la norma de la proyección de los datos sobre el
complemento ortogonal del subespacio generado por las exponenciales. La
dimensión del problema de optimización se reduce de $M(P+1)$ a $M$
variables complejas, y la estimación resultante de
$\boldsymbol{\lambda}$ suprime el sesgo de los métodos de regresión lineal
\cite{askham2018,brunton2022}. En este trabajo, los polos obtenidos por el
algoritmo Multi-TK se emplean como inicialización de
\eqref{eq:varpro_projected}, combinando la robustez del subespacio de señal
de la SVD con la consistencia estadística del ajuste no lineal.

\subsection{Discriminación de Modos Dominantes por Energía Modal}

Los estimadores con orden sobredimensionado (Multi-TK con $L>M$, VarPro
con modos triviales adicionales para absorber ruido) entregan un conjunto
de modos en el que los modos físicos dominantes están mezclados con modos
triviales. Para separarlos, además del criterio sobre valores singulares
\eqref{eq:energy}, se emplea el criterio de \textbf{energía modal} sobre
las amplitudes estimadas \cite{jiang2021}: la energía del modo $k$ se
define a partir de las amplitudes complejas de todos los canales,
\begin{equation}
    E_k \;=\; \sum_{i=1}^{P} |a_{i,k}|^2
    \;=\; \lVert \mathbf{s}_k \rVert_2^2,
    \label{eq:mode_energy}
\end{equation}
y los modos se ordenan por energía decreciente. El peso de energía
acumulada
\begin{equation}
    E_{s,k} \;=\;
    \frac{\displaystyle\sum_{i=1}^{k} E_i}{\displaystyle\sum_{i=1}^{M} E_i}
    \times 100\%
    \label{eq:cum_mode_energy}
\end{equation}
determina los $k$ primeros modos como dominantes cuando
$E_{s,k} \ge E_0$, con $E_0 = 85\%$ como valor validado empíricamente
frente al análisis de eigenvalores y a la densidad espectral de potencia
en el sistema de prueba de 16 generadores y 68 buses \cite{jiang2021}. A
diferencia de \eqref{eq:energy}, que opera sobre el subespacio de la
matriz de datos \emph{antes} de estimar los polos, el criterio
\eqref{eq:cum_mode_energy} opera sobre la contribución de cada modo ya
estimado a la reconstrucción de las señales, por lo que ambos criterios
son complementarios: el primero fija el orden del modelo $M$ y el segundo
jerarquiza los modos físicos dentro de ese orden.

\subsection{Forma Modal (Mode Shape)}

La \textbf{forma modal} del modo $k$ describe la distribución espacial:
\begin{equation}
\mathbf{s}_k =
\begin{bmatrix}
    a_{1,k} \\
    a_{2,k} \\
    \vdots \\
    a_{P,k}
\end{bmatrix}
=
\begin{bmatrix}
    |a_{1,k}|\,e^{j\phi_{1,k}} \\
    |a_{2,k}|\,e^{j\phi_{2,k}} \\
    \vdots \\
    |a_{P,k}|\,e^{j\phi_{P,k}}
\end{bmatrix}
\end{equation}

La magnitud $|\mathbf{s}_k|$ indica observabilidad; la fase $\angle\mathbf{s}_k$ revela relaciones entre señales. Para comparar formas
modales entre modos o entre métodos, $\mathbf{s}_k$ se normaliza a norma
unitaria, $\hat{\mathbf{s}}_k = \mathbf{s}_k / \lVert\mathbf{s}_k\rVert_2$,
o respecto a su componente de mayor magnitud, convención empleada en la
literatura de identificación basada en mediciones
\cite{jiang2021,li2020era}.

La forma modal empírica $\mathbf{s}_k$ tiene una interpretación exacta en
términos del modelo de estados \eqref{eq:statespace}: por
\eqref{eq:measured_modal}, $a_{i,k} =
(\mathbf{c}_i^\top\boldsymbol{\phi}_k)\,c_k$, de modo que
\begin{equation}
    \mathbf{s}_k \;=\; c_k
    \begin{bmatrix}
        \mathbf{c}_1^\top \boldsymbol{\phi}_k \\
        \vdots \\
        \mathbf{c}_P^\top \boldsymbol{\phi}_k
    \end{bmatrix},
    \label{eq:shape_link}
\end{equation}
es decir, $\mathbf{s}_k$ es la proyección del \textbf{eigenvector derecho}
$\boldsymbol{\phi}_k$ sobre los canales medidos, escalada por el factor de
excitación $c_k$, común a todos los canales. Dado que $c_k$ es un escalar,
las magnitudes \textit{relativas} y las fases \textit{relativas} entre
canales de $\mathbf{s}_k$ coinciden con las del eigenvector derecho, que es
precisamente la forma modal del análisis de pequeña señal
\cite{kundur1994}. En particular, los canales en contrafase
($\Delta\phi \approx 180^\circ$) identifican los grupos coherentes de
generadores que oscilan uno contra otro en un modo interárea.

\subsection{Factor de Participación}

Para el modo $k$, la participación normalizada de la señal $i$ es:
\begin{equation}
\hat{p}_{ik} = \frac{|a_{i,k}|}{\displaystyle\sum_{j=1}^{P} |a_{j,k}|}
\end{equation}

Debe distinguirse $\hat{p}_{ik}$ del \textbf{factor de participación} del
análisis de pequeña señal, definido como el producto de los elementos de
los eigenvectores derecho e izquierdo,
\begin{equation}
    p_{ki} \;=\; \phi_{ki}\,\psi_{ik},
    \label{eq:participation_state}
\end{equation}
que mide la participación neta del estado $k$ en el modo $i$ combinando
actividad (eigenvector derecho) y controlabilidad de la condición inicial
(eigenvector izquierdo), y cuya evaluación exige el conocimiento explícito
de la matriz $\mathbf{A}$ \cite{kundur1994}. El índice $\hat{p}_{ik}$ aquí
propuesto es, en cambio, una medida \textit{empírica de observabilidad
relativa} del modo $k$ en el canal $i$, calculable exclusivamente a partir
de mediciones sincrofasoriales sin requerir un modelo del sistema. Ambas
nociones coinciden en su uso práctico: jerarquizar las señales (y por ende
las máquinas o áreas) dominantes en cada modo e identificar los grupos
coherentes de generadores, lo que habilita la localización de los canales
óptimos para monitoreo y control de oscilaciones \cite{sauer2017}. Esta
distinción está respaldada por la literatura de identificación basada en
mediciones: los métodos que estiman explícitamente una realización de
estado (p.~ej.\ ERA) recuperan el factor de participación formal a partir
de los eigenvectores derecho e izquierdo de la matriz de estados
realizada \cite{li2020era}, mientras que los métodos de ajuste de
amplitudes (OptVarPro, Prony multicanal, Multi-TK) construyen índices de
participación a partir de las magnitudes $|a_{i,k}|$
\cite{jiang2021}. En ambos enfoques se ha verificado, sobre el sistema de
prueba de 16 generadores y 68 buses, que los generadores con mayor índice
empírico coinciden con los de mayor factor de participación obtenido por
análisis de eigenvalores, aun cuando los valores numéricos difieren
\cite{jiang2021,li2020era}.

\subsection{Identificación de Grupos Coherentes: Coseno Director y
Agrupamiento Espectral}

La forma modal de un solo modo separa a los generadores en, a lo sumo,
dos grupos en contrafase; la partición del sistema en grupos coherentes
bajo la acción conjunta de \emph{múltiples} modos dominantes requiere una
medida de similitud entre generadores construida sobre el conjunto
completo de formas modales \cite{jiang2021,li2018cwt}. Sea
$\mathbf{r}_i = [a_{i,1}, a_{i,2}, \ldots, a_{i,M}] \in \mathbb{C}^{1\times M}$
la fila de amplitudes complejas (residuos) del generador $i$ sobre los $M$
modos dominantes, normalizada a norma unitaria,
$\hat{\mathbf{r}}_i = \mathbf{r}_i / \lVert\mathbf{r}_i\rVert_2$. El
\textbf{coseno director} entre los generadores $i$ y $j$ se define como la
parte real del producto interno complejo normalizado:
\begin{equation}
    r_{ij} \;=\;
    \operatorname{Re}\bigl( \hat{\mathbf{r}}_i\, \hat{\mathbf{r}}_j^{H} \bigr)
    \;\in\; [-1,\, 1].
    \label{eq:direction_cosine}
\end{equation}
Dos generadores completamente coherentes (misma respuesta modal, fase
relativa nula) tienen $r_{ij}=1$; dos generadores en contrafase perfecta
tienen $r_{ij}=-1$ \cite{li2018cwt,li2020era}. Es esencial conservar el
signo en \eqref{eq:direction_cosine}: tomar la magnitud
$|\hat{\mathbf{r}}_i \hat{\mathbf{r}}_j^{H}|$ colapsa el rango a $[0,1]$ y
asigna similitud máxima a generadores en oposición de fase, que por
definición pertenecen a grupos coherentes distintos.

A partir de la matriz de cosenos directores
$\mathbf{R} = [r_{ij}] \in \mathbb{R}^{n_g \times n_g}$ existen dos
estrategias de partición:

\textbf{(a) Umbralización.} Dos generadores se declaran coherentes si
$r_{ij} \ge r_0$, con $r_0 = 0.8$ como umbral validado en la literatura
\cite{li2018cwt,li2020era}. La matriz binaria de adyacencia resultante
define un grafo no dirigido cuyas componentes conexas son los grupos
coherentes. Su limitación es la necesidad de fijar $r_0$ a priori.

\textbf{(b) Agrupamiento espectral con eigengap.} Para evitar el umbral,
se construye un grafo ponderado $G=(V,E)$ con vértices en los generadores
y pesos de arista \cite{jiang2021}
\begin{equation}
    w_{ij} \;=\;
    \begin{cases}
        r_{ij}, & r_{ij} \ge 0,\\
        0, & r_{ij} < 0,
    \end{cases}
    \qquad i \ne j,
    \label{eq:edge_weights}
\end{equation}
de modo que la incoherencia (signo negativo) se traduce en ausencia de
arista. Con $d_i = \sum_{j} w_{ij}$ y
$\mathbf{D} = \operatorname{diag}(d_1,\ldots,d_{n_g})$, el laplaciano del
grafo y su versión normalizada simétrica son
\begin{equation}
    \mathbf{L} \;=\; \mathbf{D} - \mathbf{W},
    \qquad
    \mathbf{L}_N \;=\; \mathbf{D}^{-1/2}\,\mathbf{L}\,\mathbf{D}^{-1/2}.
    \label{eq:laplacian}
\end{equation}
Los eigenvalores de $\mathbf{L}_N$, ordenados ascendentemente
$0 = \mu_1 \le \mu_2 \le \cdots \le \mu_{n_g}$, codifican la estructura de
comunidades del grafo: el número de grupos coherentes $K$ se selecciona
mediante el criterio del \textbf{eigengap},
\begin{equation}
    K \;=\; \operatorname*{arg\,max}_{k}\;
    \bigl( \mu_{k+1} - \mu_k \bigr),
    \label{eq:eigengap}
\end{equation}
es decir, $K$ es el índice donde ocurre la mayor brecha espectral
\cite{jiang2021}. La partición se obtiene aplicando $k$-medias sobre las
filas, normalizadas a norma unitaria, de la matriz formada por los $K$
primeros eigenvectores de $\mathbf{L}_N$. Este procedimiento resuelve de
forma aproximada el problema de corte de grafo y no requiere conocer
$K$ ni $r_0$ a priori; sobre el sistema de 16 generadores y 68 buses
reproduce los grupos coherentes del análisis de eigenvalores, con la
capacidad adicional de aislar generadores de comportamiento singular en
grupos unitarios \cite{jiang2021}.

\chapterbibliography